\documentstyle{article}
\begin{document}

\centerline{\bf Nonlinear Differential Equations with Inexactly Known}
\centerline{\bf Boundary Conditions and Gr\"{o}bner bases}

\vskip 0.3cm

\centerline{Vladimir P. Gerdt}
\centerline{Laboratory of Computing Techniques and Automation}
\centerline{Joint Institute for Nuclear Research}
\centerline{141980 Dubna, Russia}
\centerline{Email: {\em gerdt@jinr.dubna.su}}
\centerline{and}
\centerline{Michael G.Dmitriev and Marina V.Nesterova}
\centerline{Program System Institute of Russian Academy of Sciences}
\centerline{152140 Pereslavl-Zalessky, Russia}

\vskip 0.3cm

An algorithm is presented for constructing approximate solutions of boundary
problems for second-order differential equations of the form
\begin{equation}
\Phi(y'',y',x)=0, \quad y(a)=y_a,\quad y_b=y_b\,, \label{nbp}
\end{equation}
where $y(x)$ is a real function, $\Phi$ is a
polynomial in its arguments over rationals, and the numbers $a,b,y_a,y_b$
are also rational. Here $y_a$ or $y_b$ or both of them are assumed to be
known with some reasonably small errors. The algorithm is based on the use
of the quadratic penalty functions for the approximately given boundary
conditions and solving the corresponding unconditional extremum problem.
In particular, if equation $\Phi(y'',y',x)=0$ admits the variational
formulation $(\phi_{y'})'-\phi_{y}=0$, and both $y_a,y_b$ are inexactly
known, then an approximate solution of~(\ref{nbp}) can be constructed
from the condition
$$
F(\{c_i\},\alpha,\beta)=\int_a^b \phi(\tilde{y}'',\tilde{y}',x)dx +
 \alpha[\tilde{y}(a)-y_a]^2 + \beta[\tilde{y}(b)-y_b]^2
 \ \Longrightarrow \ \min_{\{c_i\}}\,,
$$
$$
 y\approx \tilde{y}=\sum_{i=1}^n c_i v_i(x)\,,
$$
The arising system of nonlinear algebraic equations in the coefficients
$c_i$ for some appropriate set of basic functions $v_i(x)$ is solved by
construction of a lexicographical Gr\"{o}bner basis. It is shown that the
construction of such a basis allows one to develop a perturbation scheme
in the inverse degrees of the penalty parameters $\alpha$ and $\beta$.
The related computational aspects for parametric Gr\"{o}bner bases are
considered. The proposed algorithm is illustrated by an example of the
boundary problem with the use of computer algebra system Reduce. The
obtained accuracy is analysed in comparison with some other methods
used to solve that particular boundary problem.

\end{document}
